\documentclass{article}

\usepackage{arxiv}

\usepackage[utf8]{inputenc}
\usepackage[T1]{fontenc}
\usepackage{hyperref}
\usepackage{url}
\usepackage{booktabs}
\usepackage{amsfonts}
\usepackage{nicefrac}
\usepackage{microtype}
\usepackage{amsmath}
\usepackage{graphicx}
\usepackage{float}
\usepackage{xcolor}
\usepackage{listings}
\graphicspath{ {./graficas_tp22/} }

\lstset{
    language=Python,
    basicstyle=\ttfamily\small,
    keywordstyle=\color{blue!70!black},
    commentstyle=\color{green!40!black},
    stringstyle=\color{red!60!black},
    numbers=left,
    numberstyle=\tiny,
    frame=single,
    breaklines=true,
    showstringspaces=false,
    tabsize=4
}

\title{TP 2.2 - GENERADORES DE NÚMEROS PSEUDOALEATORIOS DE DISTINTAS DISTRIBUCIONES}

\author{
  Brizio Augusto \\
  Universidad Tecnológica Nacional - FRRO \\
  Zeballos 1341, S2000, Rosario, Argentina \\
  \texttt{brizioaugusto@gmail.com} \\
  \And
  Casermeiro Gonzalo Nicolás \\
  Universidad Tecnológica Nacional - FRRO \\
  Zeballos 1341, S2000, Rosario, Argentina \\
  \texttt{gonzacasermeiro@gmail.com} \\
  \And
  Messina Tiziano Leonel \\
  Universidad Tecnológica Nacional - FRRO \\
  Zeballos 1341, S2000, Rosario, Argentina \\
  \texttt{tizilmessina@gmail.com} \\
  \And
  De la Rosa Valentin \\
  Universidad Tecnológica Nacional - FRRO \\
  Zeballos 1341, S2000, Rosario, Argentina \\
  \texttt{mferreyra079@gmail.com} \\
}

\begin{document}
\maketitle

\begin{abstract}
Este trabajo continúa el estudio de generadores pseudoaleatorios iniciado en el
TP 2.1. A partir de números uniformes en el intervalo $[0,1)$ se implementaron
generadores para nueve distribuciones de probabilidad: uniforme, exponencial,
gamma, normal, Pascal, binomial, hipergeométrica, Poisson y empírica discreta.
El objetivo fue transformar una fuente uniforme en variables aleatorias con
comportamientos probabilísticos distintos, utilizando transformación inversa,
composición de variables, ensayos de Bernoulli, muestreo sin reemplazo y el
método de Box-Muller.

Para validar las simulaciones se generaron $n=10\,000$ valores por distribución
y se compararon la media y la varianza observadas contra sus valores teóricos.
Además, para distribuciones continuas se aplicó una prueba de
Kolmogorov-Smirnov y para distribuciones discretas una medida de ajuste
chi-cuadrado. Los resultados muestran errores relativos bajos en media y
varianza, lo que confirma que las transformaciones implementadas producen
secuencias compatibles con las distribuciones esperadas.
\end{abstract}

\section{Introducción}

En el TP 2.1 se estudió la generación de números pseudoaleatorios uniformes y
la necesidad de verificar su calidad estadística. Sin embargo, muchas
simulaciones no trabajan directamente con valores uniformes, sino con tiempos
entre llegadas, cantidades de eventos, demandas, tiempos de servicio, fallas,
éxitos o variables físicas que responden a otras distribuciones.

El problema central es entonces transformar una secuencia uniforme $U(0,1)$ en
observaciones de una distribución deseada. Naylor presenta este enfoque como la
base para generar valores de variables aleatorias en simulación computacional:
partir de números uniformes y aplicar funciones matemáticas o algoritmos
específicos para obtener la distribución buscada \cite{naylor}.

El presente informe implementa y evalúa generadores para distribuciones
continuas y discretas. En cada caso se presenta la función de probabilidad o
densidad, sus parámetros, media y varianza, y el método computacional utilizado.
El programa fue escrito en Python 3.x y genera automáticamente gráficas y una
tabla resumen para contrastar el comportamiento simulado con el esperado.

\section{Métodos}

Todos los generadores parten de números pseudoaleatorios uniformes $U \sim
U(0,1)$ provistos por \texttt{random.random()}. Se utilizó una semilla fija
(\texttt{42}) para reproducir exactamente los resultados y $n=10\,000$ muestras
por distribución.

\subsection{Distribución uniforme}

La distribución uniforme continua en $(a,b)$ posee densidad:

\begin{equation}
    f(x) =
    \begin{cases}
        \frac{1}{b-a}, & a \leq x \leq b \\
        0, & \text{en otro caso}
    \end{cases}
\end{equation}

\noindent Su acumulada es:

\begin{equation}
    F(x) = \frac{x-a}{b-a}
\end{equation}

\noindent Aplicando transformación inversa:

\begin{equation}
    U = \frac{x-a}{b-a}
    \quad \Rightarrow \quad
    X = a + (b-a)U
\end{equation}

Su media y varianza son:

\begin{equation}
    E[X] = \frac{a+b}{2}, \qquad Var(X) = \frac{(b-a)^2}{12}
\end{equation}

En la simulación se utilizó $a=2$ y $b=8$.

\subsection{Distribución exponencial}

La distribución exponencial modela tiempos entre eventos independientes con
tasa $\lambda$:

\begin{equation}
    f(x)=\lambda e^{-\lambda x}, \qquad x \geq 0
\end{equation}

Su función acumulada es:

\begin{equation}
    F(x)=1-e^{-\lambda x}
\end{equation}

La inversa se obtiene igualando $U=F(x)$:

\begin{equation}
    U = 1-e^{-\lambda x}
    \quad \Rightarrow \quad
    X = -\frac{\ln(1-U)}{\lambda}
\end{equation}

Como $1-U$ también es uniforme, suele escribirse $X=-\ln(U)/\lambda$. En este
trabajo se usó $\lambda=1{,}5$.

\subsection{Distribución gamma}

La distribución gamma, en la parametrización usada por Naylor, puede
interpretarse como la suma de $k$ variables exponenciales independientes con la
misma tasa $\lambda$ cuando $k$ es entero positivo. Su densidad es:

\begin{equation}
    f(x)=\frac{\lambda^k x^{k-1}e^{-\lambda x}}{(k-1)!}, \qquad x \geq 0
\end{equation}

\noindent con:

\begin{equation}
    E[X]=\frac{k}{\lambda}, \qquad Var(X)=\frac{k}{\lambda^2}
\end{equation}

Para $k$ entero, se puede generar una variable gamma/Erlang sumando
exponenciales:

\begin{equation}
    X = \sum_{i=1}^{k} -\frac{\ln(1-U_i)}{\lambda}
\end{equation}

En la simulación se empleó $k=3$ y $\lambda=1{,}2$.

\subsection{Distribución normal}

La distribución normal con media $\mu$ y desvío estándar $\sigma$ tiene densidad:

\begin{equation}
    f(x)=\frac{1}{\sigma\sqrt{2\pi}}
    e^{-\frac{1}{2}\left(\frac{x-\mu}{\sigma}\right)^2}
\end{equation}

Para generarla se utilizó el método directo de Box-Muller. Dados dos uniformes
independientes $U_1$ y $U_2$:

\begin{equation}
    Z_1=\sqrt{-2\ln(U_1)}\cos(2\pi U_2), \qquad
    Z_2=\sqrt{-2\ln(U_1)}\sin(2\pi U_2)
\end{equation}

\noindent producen dos normales estándar. Luego se transforma:

\begin{equation}
    X = \mu + \sigma Z
\end{equation}

Se utilizaron $\mu=10$ y $\sigma=2$.

\subsection{Distribución de Pascal}

La distribución de Pascal o binomial negativa modela la cantidad de fallas $X$
antes de obtener $r$ éxitos, con probabilidad de éxito $p$:

\begin{equation}
    P(X=x)=\binom{r+x-1}{x}p^r(1-p)^x,
    \qquad x=0,1,2,\ldots
\end{equation}

\begin{equation}
    E[X]=\frac{r(1-p)}{p}, \qquad
    Var(X)=\frac{r(1-p)}{p^2}
\end{equation}

El algoritmo simula ensayos de Bernoulli hasta acumular $r$ éxitos. En el
experimento se usó $r=4$ y $p=0{,}35$.

\subsection{Distribución binomial}

La distribución binomial cuenta la cantidad de éxitos en $n$ ensayos
independientes de Bernoulli con probabilidad $p$:

\begin{equation}
    P(X=x)=\binom{n}{x}p^x(1-p)^{n-x}
\end{equation}

\begin{equation}
    E[X]=np, \qquad Var(X)=np(1-p)
\end{equation}

Se generó realizando $n=12$ ensayos por observación y contando los uniformes
menores a $p=0{,}4$.

\subsection{Distribución hipergeométrica}

La distribución hipergeométrica representa extracciones sin reemplazo desde una
población de tamaño $N$, con $K$ elementos exitosos y $n$ extracciones
realizadas:

\begin{equation}
    P(X=x)=
    \frac{\binom{K}{x}\binom{N-K}{n-x}}{\binom{N}{n}}
\end{equation}

\begin{equation}
    E[X]=n\frac{K}{N}, \qquad
    Var(X)=n\frac{K}{N}\left(1-\frac{K}{N}\right)\frac{N-n}{N-1}
\end{equation}

El algoritmo actualiza las cantidades restantes de éxitos y fracasos después de
cada extracción. Se usó $N=50$, $K=18$ y $n=10$.

\subsection{Distribución de Poisson}

La distribución de Poisson modela la cantidad de eventos en un intervalo fijo
cuando estos ocurren con tasa promedio $\lambda$:

\begin{equation}
    P(X=x)=\frac{e^{-\lambda}\lambda^x}{x!},
    \qquad x=0,1,2,\ldots
\end{equation}

\begin{equation}
    E[X]=\lambda, \qquad Var(X)=\lambda
\end{equation}

Se utilizó el método multiplicativo: se generan uniformes y se acumula su
producto hasta que cae por debajo de $e^{-\lambda}$. El número de factores
menos uno es el valor de Poisson. Se usó $\lambda=4$.

\subsection{Distribución empírica discreta}

Una distribución empírica discreta asigna probabilidades explícitas a un
conjunto finito de valores. En este trabajo:

\begin{center}
\begin{tabular}{cccccc}
\toprule
$x$ & 1 & 2 & 3 & 4 & 5 \\
\midrule
$P(X=x)$ & 0{,}10 & 0{,}20 & 0{,}35 & 0{,}25 & 0{,}10 \\
\bottomrule
\end{tabular}
\end{center}

El método consiste en formar probabilidades acumuladas y seleccionar el primer
valor cuyo acumulado supere al uniforme generado:

\begin{equation}
    X=x_i \quad \text{si} \quad F(x_{i-1}) < U \leq F(x_i)
\end{equation}

\section{Implementación en Python}

El Código~\ref{cod:continuas} muestra las transformaciones principales para
distribuciones continuas. Todas parten de valores uniformes generados por
\texttt{random.random()}.

\begin{lstlisting}[caption={Generadores continuos implementados en Python.}, label={cod:continuas}]
def uniforme(n, a=2.0, b=8.0):
    return [a + (b - a) * random.random() for _ in range(n)]

def exponencial(n, lambd=1.5):
    return [-math.log(1 - random.random()) / lambd for _ in range(n)]

def gamma_erlang(n, k=3, lambd=1.2):
    valores = []
    for _ in range(n):
        total = 0.0
        for _ in range(k):
            total += -math.log(1 - random.random()) / lambd
        valores.append(total)
    return valores

def normal_box_muller(n, media=10.0, desvio=2.0):
    valores = []
    while len(valores) < n:
        u1 = max(random.random(), 1e-12)
        u2 = random.random()
        radio = math.sqrt(-2 * math.log(u1))
        z1 = radio * math.cos(2 * math.pi * u2)
        z2 = radio * math.sin(2 * math.pi * u2)
        valores.append(media + desvio * z1)
        if len(valores) < n:
            valores.append(media + desvio * z2)
    return valores
\end{lstlisting}

El Código~\ref{cod:discretas} resume la lógica de generación para algunas
distribuciones discretas. En binomial y Pascal se emplean ensayos de Bernoulli;
en Poisson, el método multiplicativo.

\begin{lstlisting}[caption={Generadores discretos implementados en Python.}, label={cod:discretas}]
def pascal(n, r=4, p=0.35):
    valores = []
    for _ in range(n):
        exitos = 0
        fallas = 0
        while exitos < r:
            if random.random() < p:
                exitos += 1
            else:
                fallas += 1
        valores.append(fallas)
    return valores

def binomial(n, ensayos=12, p=0.4):
    valores = []
    for _ in range(n):
        exitos = sum(1 for _ in range(ensayos) if random.random() < p)
        valores.append(exitos)
    return valores

def poisson(n, lambd=4.0):
    valores = []
    limite = math.exp(-lambd)
    for _ in range(n):
        k = 0
        producto = 1.0
        while producto > limite:
            k += 1
            producto *= random.random()
        valores.append(k - 1)
    return valores
\end{lstlisting}

\section{Resultados y discusión}

La Tabla~\ref{tab:resultados} presenta la comparación entre valores observados
y teóricos para $10\,000$ muestras. En distribuciones continuas, el último
estadístico corresponde a Kolmogorov-Smirnov; en distribuciones discretas,
corresponde a chi-cuadrado.

\begin{table}[H]
\centering
\caption{Resumen de resultados de la simulación con semilla 42.}
\label{tab:resultados}
\begin{tabular}{lrrrrr}
\toprule
\textbf{Distribución} & \textbf{Media obs.} & \textbf{Media teo.} &
\textbf{Var. obs.} & \textbf{Var. teo.} & \textbf{KS/$\chi^2$} \\
\midrule
Uniforme & 5{,}00151 & 5{,}00000 & 2{,}97963 & 3{,}00000 & 0{,}00626 \\
Exponencial & 0{,}67046 & 0{,}66667 & 0{,}43472 & 0{,}44444 & 0{,}00984 \\
Gamma & 2{,}48818 & 2{,}50000 & 2{,}06080 & 2{,}08333 & 0{,}01039 \\
Normal & 10{,}02792 & 10{,}00000 & 3{,}99574 & 4{,}00000 & 0{,}01046 \\
Pascal & 7{,}42330 & 7{,}42857 & 20{,}91281 & 21{,}22449 & 31{,}18428 \\
Binomial & 4{,}77190 & 4{,}80000 & 2{,}83935 & 2{,}88000 & 8{,}76653 \\
Hipergeométrica & 3{,}60420 & 3{,}60000 & 1{,}90193 & 1{,}88082 & 4{,}72206 \\
Poisson & 4{,}00470 & 4{,}00000 & 3{,}97608 & 4{,}00000 & 13{,}13614 \\
Empírica discreta & 3{,}05630 & 3{,}05000 & 1{,}24505 & 1{,}24750 & 2{,}57414 \\
\bottomrule
\end{tabular}
\end{table}

Los errores relativos fueron bajos en todos los casos. La media observada de la
uniforme fue $5{,}00151$, prácticamente igual al valor teórico $5$. La
exponencial y la gamma muestran pequeñas diferencias atribuibles al error
muestral. La normal generada por Box-Muller también coincide de forma clara con
la media y varianza esperadas.

En las distribuciones discretas, la comparación mediante barras simuladas y
probabilidades teóricas permite verificar visualmente el ajuste. La binomial,
hipergeométrica, Poisson y empírica discreta se comportan de manera muy estable.
La distribución de Pascal presenta mayor dispersión, algo esperable por su
varianza elevada.

\subsection{Gráficas de distribuciones continuas}

\begin{figure}[H]
    \centering
    \includegraphics[width=\linewidth]{01_uniforme.png}
    \caption{Valores generados para una distribución uniforme $U(2,8)$.}
    \label{fig:uniforme}
\end{figure}

\begin{figure}[H]
    \centering
    \includegraphics[width=\linewidth]{02_exponencial.png}
    \caption{Valores generados para una distribución exponencial con $\lambda=1{,}5$.}
    \label{fig:exponencial}
\end{figure}

\begin{figure}[H]
    \centering
    \includegraphics[width=\linewidth]{03_gamma.png}
    \caption{Valores generados para una distribución gamma/Erlang con $k=3$ y $\lambda=1{,}20$.}
    \label{fig:gamma}
\end{figure}

\begin{figure}[H]
    \centering
    \includegraphics[width=\linewidth]{04_normal.png}
    \caption{Valores generados para una distribución normal $N(10,2)$.}
    \label{fig:normal}
\end{figure}

\subsection{Gráficas de distribuciones discretas}

\begin{figure}[H]
    \centering
    \includegraphics[width=\linewidth]{05_pascal.png}
    \caption{Valores generados para una distribución de Pascal con $r=4$ y $p=0{,}35$.}
    \label{fig:pascal}
\end{figure}

\begin{figure}[H]
    \centering
    \includegraphics[width=\linewidth]{06_binomial.png}
    \caption{Valores generados para una distribución binomial con $n=12$ y $p=0{,}40$.}
    \label{fig:binomial}
\end{figure}

\begin{figure}[H]
    \centering
    \includegraphics[width=\linewidth]{07_hipergeometrica.png}
    \caption{Valores generados para una distribución hipergeométrica con $N=50$, $K=18$ y $n=10$.}
    \label{fig:hipergeometrica}
\end{figure}

\begin{figure}[H]
    \centering
    \includegraphics[width=\linewidth]{08_poisson.png}
    \caption{Valores generados para una distribución de Poisson con $\lambda=4$.}
    \label{fig:poisson}
\end{figure}

\begin{figure}[H]
    \centering
    \includegraphics[width=\linewidth]{09_empirica_discreta.png}
    \caption{Valores generados para una distribución empírica discreta.}
    \label{fig:empirica}
\end{figure}

\subsection{Comparación global}

\begin{figure}[H]
    \centering
    \includegraphics[width=\linewidth]{10_comparacion_medias.png}
    \caption{Comparación entre medias observadas y medias teóricas.}
    \label{fig:medias}
\end{figure}

\begin{figure}[H]
    \centering
    \includegraphics[width=\linewidth]{11_comparacion_varianzas.png}
    \caption{Comparación entre varianzas observadas y varianzas teóricas.}
    \label{fig:varianzas}
\end{figure}

\section{Conclusiones}

El trabajo demuestra que una fuente uniforme validada puede transformarse en
variables aleatorias de múltiples distribuciones. La transformación inversa
resultó directa para la uniforme y la exponencial; la gamma se generó como suma
de exponenciales; la normal mediante Box-Muller; y las distribuciones discretas
mediante ensayos, acumulación de probabilidades o mecanismos propios del
experimento probabilístico.

Los resultados obtenidos fueron consistentes con la teoría. Las medias y
varianzas simuladas se mantuvieron cercanas a sus valores esperados, y las
gráficas muestran una buena coincidencia entre las frecuencias observadas y las
funciones teóricas. Esto confirma que los algoritmos implementados son adecuados
para simulaciones académicas y permiten representar fenómenos más generales que
los valores uniformes originales.

También se observa que cada distribución requiere una técnica de generación
distinta. No existe un único procedimiento universal igualmente simple para
todos los casos: algunas distribuciones surgen por inversión de la acumulada,
otras por suma o composición de variables, y otras por simular el proceso físico
o lógico que les da origen.

\bibliographystyle{unsrt}

\begin{thebibliography}{9}

\bibitem{naylor}
T.~H. Naylor, \textit{Técnicas de simulación en computadoras},
Editorial Limusa, 1982.

\bibitem{torres_tp22}
J.~I. Torres, \textit{TP 2.2 - Generadores de números pseudoaleatorios de
distintas distribuciones de probabilidad}, Universidad Tecnológica Nacional -
FRRO, 2026.

\bibitem{python_random}
Python Software Foundation, \textit{random --- Generate pseudo-random numbers},
\url{https://docs.python.org/3/library/random.html}.

\bibitem{box_muller}
G.~E.~P. Box and M.~E. Muller, ``A Note on the Generation of Random Normal
Deviates'', \textit{The Annals of Mathematical Statistics}, vol.~29, no.~2,
pp.~610--611, 1958.

\end{thebibliography}

\end{document}
